\section{Seis folhas, dois comprimentos de caminho}
\label{sec:bloco-seis-folhas}
\index{arvore de merkle@árvore de Merkle!prova de inclusao@prova de inclusão}

Temos já todos os dados para responder a uma pergunta concreta. Alguém
conservou a transacção cujo identificador começa por \texttt{447c77a6}.
Como lhe demonstrar que esse identificador figura, naquela posição,
no bloco \(1\,582\,196\)? Comecemos por reconstruir a árvore inteira.
Veremos depois quanto dela basta entregar a esse leitor.

Escrevamos \(h_0\) para o identificador da transacção de mineiro e
\(h_1,\ldots,h_5\) para os cinco identificadores de \texttt{tx\_hashes},
pela ordem impressa. Na tabela do bloco, a nossa transacção figura
na quarta posição do vector. Na árvore ocupa a posição 4, contando
de zero e contando a transacção de mineiro. Para dois resumos \(a,b\)
de 32 bytes escrevamos \([a,b]=K(a\|b)\), onde \(K\) designa
Keccak--256. O resumo da esquerda ocupa os primeiros 32 bytes.
O da direita ocupa os 32 seguintes. Cada aplicação de \(K\) a estes
64 bytes restitui outros 32.

No nível imediatamente acima das folhas mais fundas há quatro lugares.
Os dois primeiros recebem directamente \(h_0,h_1\). Os restantes
recebem \(U=[h_2,h_3]\) e \(V=[h_4,h_5]\). São estes quatro objectos,
e já todos do mesmo comprimento, que entram na árvore regular de
quatro lugares. Formam-se \(A=[h_0,h_1]\), \(B=[U,V]\) e finalmente
\(R=[A,B]\). A disposição mostra por que motivo há folhas a duas
distâncias diferentes da raiz.

\begin{center}
\begin{forest}
 for tree={math content, l sep=6mm, s sep=7mm, inner sep=2pt},
 [R [A [h_0] [h_1]] [B [U [h_2] [h_3]] [V [h_4] [h_5]]]]
\end{forest}
\end{center}

Cada par na lista seguinte dá, em duas metades consecutivas, a
codificação de um único resultado. A primeira linha pode ser refeita
a partir dos identificadores que começam por \texttt{6d31593c} e
\texttt{329e9c0c}. A segunda usa os que começam por
\texttt{447c77a6} e \texttt{90a698b0}.

\begin{center}\small
\begin{tabular}{c l}
Ponto da árvore & Resumo de 32 bytes\\ \hline
\(U\) & \texttt{fa0b6468cee78e59eac0cacbbb263551}\\
 & \texttt{caff0725a93ff2d090c4188a57039f0b}\\
\(V\) & \texttt{d64da9a9af3d3255122e4357cd7fdf2f}\\
 & \texttt{e7b923df7951424e7d5b92a5168fa271}\\
\(A\) & \texttt{3709091d920b22a92832da2f86866526}\\
 & \texttt{dc987dbf2b18b9121b61e7ab5081b65a}\\
\(B\) & \texttt{d2b3ae5b5e066333c97e9bd42667e40e}\\
 & \texttt{1ca7e1e9288b76144b873a9f5a9efb4a}\\
\(R\) & \texttt{605a61a5619b3fde0deaad0c79205a0d}\\
 & \texttt{8b92141acfcd04031b0d4118dc2e26cb}
\end{tabular}
\end{center}

Voltemos ao leitor que conserva \(h_4\). Entreguemos-lhe, pela
ordem em que lhe vão fazer falta, \(h_5,U,A\). Conhecendo o número
de folhas e a posição 4, o leitor conhece também o lado em que entra
cada irmão. Primeiro calcula \([h_4,h_5]=V\). Agora vem da direita e
calcula \([U,V]=B\). Ainda da direita, calcula \([A,B]=R\). Foram-lhe
dados três resumos, 96 bytes ao todo. Para outra pessoa, interessada
no primeiro pagamento ordinário, a resposta é mais curta. Na posição 1
chegam \(h_0\) e \(B\). O cálculo começa com \([h_0,h_1]\) e acaba com
\([A,B]\). Pouparam-se 32 bytes porque o primeiro pagamento já entrava
directamente no nível dos quatro lugares.

A mesma conta resolve o caso geral. Para \(n\geq 2\) folhas, seja
\(m=2^k\) a maior potência de dois que cabe em \(n\). Procuremos, no
nível de \(m\) lugares, o número \(s\) de folhas que entram sozinhas
e o número \(p\) de pares. Cada lugar recebe um objecto, donde
\(s+p=m\). Cada par consome duas folhas, donde \(s+2p=n\). Subtraindo
as igualdades ficamos com
\[
                    p=n-m,\qquad s=2m-n .
\]
Entram primeiro as \(s\) folhas directas, com caminhos de \(k\)
irmãos. As \(2p\) restantes têm caminhos de \(k+1\) irmãos. Para as
seis folhas do bloco, \(m=4\), \(s=2\) e \(p=2\). Para oito folhas
resultam oito caminhos com três irmãos. E se houver uma única
transacção, a de mineiro, o seu identificador já é a raiz. Uma prova
de inclusão declara assim a folha, a sua posição e o número de
folhas relativamente aos quais se há-de interpretar a sucessão de
irmãos \cite{monero-source-v16}.

\subsection{Como se calcula a folha do mineiro}

Na construção anterior aceitámos o identificador \(h_0\) que acompanha
a descrição do bloco. Podemos agora obtê-lo dos campos de
\texttt{miner\_tx}. Esta é a única folha cujo conteúdo se encontra
por inteiro no próprio bloco. Na primeira leitura, os cinco
identificadores ordinários vieram prontos a usar. A transacção do
mineiro chegou com uma entrada de geração, uma saída, a sua chave
de destino e os dados adicionais. É a partir desses elementos que o
nó forma a primeira folha.

O byte \texttt{02} dá a versão da transacção. Seguem-se os três bytes
\texttt{b0 c9 60}, que exprimem a altura de desbloqueio
\(48+73\cdot128+96\cdot128^2=1\,582\,256\).
O byte \texttt{01} conta uma entrada e \texttt{ff} identifica a
geração. Segue-se a altura do próprio bloco, \texttt{f4 c8 60},
ou \(116+72\cdot128+96\cdot128^2=1\,582\,196\).
Reconhecem-se na diferença as sessenta alturas que têm de decorrer
para a saída de mineiro poder ser gasta.

O contador de saídas vale um. O montante público ocupa sete bytes,
\texttt{97 80 b5 85 f2 86 01}. A etiqueta \texttt{02} identifica
a saída por chave. Seguem-se os 32 bytes da chave e, por último,
\texttt{2b} e os 43 bytes do campo \texttt{extra}. O comprimento
do prefixo resulta da soma
\[
 \underbrace{1+3}_{\text{versão e desbloqueio}}+
 \underbrace{1+1+3}_{\text{entrada}}+
 \underbrace{1+7+1+32}_{\text{saída}}+
 \underbrace{1+43}_{\text{dados adicionais}}=94 .
\]
O resumo \(h_p\) destes 94 bytes começa por \texttt{1412cf55}.
A base RingCT limita-se ao byte \texttt{00}, a indicação de tipo
nulo. O seu resumo \(h_b=K(\mathtt{00})\) começa por
\texttt{bc36789e}. A terceira parcela do identificador é aqui o
valor nulo de 32 bytes, \(0^{256}\). Podemos finalmente calcular
\[
                   h_0=K(h_p\|h_b\|0^{256}).
\]
O resultado começa por \texttt{06fb3e1c} e termina em
\texttt{0557e6e9}, como no identificador completo impresso acima.
Distinguem-se assim duas ocorrências de zero que ocupam lugares
diferentes na conta. O byte que declara o tipo entra na função
de resumo da base. Os 32 bytes nulos entram directamente na terceira
posição do argumento final. Regressando ao caminho de \(h_1\),
dispomos agora de uma maneira de refazer o seu primeiro irmão a
partir dos dados do próprio bloco.

\section{Da raiz ao identificador impresso}

Até aqui apurámos que \(h_4\) ocupa um lugar preciso numa árvore de
raiz \(R\). Falta prender essa raiz ao bloco cujo identificador o nó
imprimiu. Reunamos para isso os bytes do cabeçalho. As duas versões
valem sete, ocupando um byte cada uma. O instante \(1\,527\,458\,161\),
escrito por grupos de sete bits a partir da direita, tem por algarismos
\(113,82,44,88,5\). Com o bit de continuação nos quatro primeiros,
fica \texttt{f1 d2 ac d8 05}. Seguem-se os 32 bytes do identificador
anterior, começado em \texttt{30bb9b47}. O \emph{nonce} é um inteiro
de quatro bytes com o menos significativo em primeiro lugar.
Como \(2\,147\,489\,363=2^{31}+22\cdot256+83\), a sua escrita é
\texttt{53 16 00 80}. O cabeçalho tem, pois, \(2+5+32+4=43\) bytes.

Designemos por \(Z\) a concatenação deste cabeçalho, dos 32 bytes
de \(R\) e do contador \texttt{06}. O conjunto tem \(43+32+1=76\)
bytes. Para obter o identificador, a serialização de \(Z\) como
cadeia binária antepõe o comprimento em inteiro de largura variável.
O número 76 cabe num byte, \texttt{4c}. O argumento de \(K\) tem,
portanto, 77 bytes. O resultado de \(K(\mathtt{4c}\|Z)\) é
\begin{center}
\texttt{50c8e5e51453c2ab85ef99d817e16654}\\
\texttt{0b40ef5fd2ed15ebc863091ca2a04594}.
\end{center}
Reconhece-se o valor do campo \texttt{hash} impresso pelo nó. A
contagem de transacções, seis, e o comprimento do objecto, 76,
figuram em extremidades opostas do argumento e medem coisas diferentes.

A quem já dispõe de um cabeçalho e de um identificador de bloco
aceite, o caminho fornece agora uma verificação curta. Reconstitui
a raiz a partir do identificador da transacção, junta-lhe o cabeçalho
e a contagem e compara o resultado final com o identificador de bloco.
A resistência a colisões dá significado a esta igualdade. Se possuir
também a transacção inteira, pode começar um passo atrás, calculando
o identificador pelos seus próprios bytes, como no
apêndice~\ref{appendix:pagamento-tipo-seis}.

O significado de «aceite» depende de outra informação. Para situar
o bloco na lista que o nó escolheu, entram as ligações aos antecessores,
a validade dos blocos e o trabalho acumulado. Para determinar a
profundidade entram ainda os sucessores conhecidos do nó nesse
instante. Nenhum desses sucessores foi necessário para refazer os
dois percursos acima. O campo \texttt{depth} pode assim passar de
dois para três quando chega outro bloco, e o leitor continuará a
usar os mesmos três irmãos para reconhecer o quarto pagamento.
